\begin{table}[tbp]
\caption{Constructs and their empirical anchors.}
\label{tab:ch2-constructs}
\begingroup\fontsize{9.2}{11.8}\selectfont
\renewcommand{\arraystretch}{1.12}
\begin{tabularx}{\textwidth}{@{}>{\raggedright\arraybackslash}p{20mm}>{\raggedright\arraybackslash}X >{\raggedright\arraybackslash}X@{}}
\toprule
\textbf{Construct} & \textbf{Evidence in our studies} & \textbf{Interpretation and boundary}\\
\midrule
Activity & Multiple behaviours, their distribution, and human/automated participation (Activity study; Xia, Ch.~5). & Recorded participation under a declared selection, weighting, and period; not a direct measure of benefit.\\[5pt]
Contribution & Coordination across issues and pull requests, including work understated by commits (ORL, Section~4.1). & Work understood in relation to a collaborative undertaking; content and outcomes require further evidence.\\[5pt]
Influence & Ecosystem and community graph constructions (OpenRank; Zhao, Chs.~5--6). & Model-dependent relational standing; graph scale and priors determine its interpretation.\\[5pt]
Value & Different repository purposes despite shared metric names (position paper, Table~4). & Stakeholder interpretation of effects relative to a reference; a ranking does not establish these effects.\\[5pt]
Sustainability & Archival explanations, alternatives, transition practices, and linked forks (Xia, Ch.~4). & Continuity of a specified function; a recorded transition does not establish successful continuation.\\
\bottomrule
\end{tabularx}
\endgroup
\tablenote{Sources: \citet{xia2022,xia2024thesis,zhao2024,zhao2025thesis,han2026rank,han2024ose}. The interpretations are this chapter's synthesis. The theses and papers belong to an overlapping research programme; they are not counted as independent replications.}
\end{table}
